\chapter{Unit 2: Computer Fundamentals}
\label{chap:unit2_qb}

\section{Practice Question Set (50 MCQs)}

\mcqitem{1}{An electronic, programmable machine that accepts data, processes it, produces output, and stores results is a/an:}{Transistor}{Computer}{Assembler}{Sensor}

\mcqitem{2}{What is the sequence of the fundamental data processing cycle in computing?}{Output $\rightarrow$ Process $\rightarrow$ Input $\rightarrow$ Storage}{Input $\rightarrow$ Process $\rightarrow$ Output $\rightarrow$ Storage}{Process $\rightarrow$ Storage $\rightarrow$ Output $\rightarrow$ Input}{Storage $\rightarrow$ Input $\rightarrow$ Process $\rightarrow$ Output}

\mcqitem{3}{Raw, unorganized facts such as numbers, names, or symbols with no immediate meaning are called:}{Information}{Data}{Knowledge}{Wisdom}

\mcqitem{4}{Processed data that has been organized to carry context, meaning, and usefulness is termed:}{Raw input}{Information}{Byte}{Firmware}

\mcqitem{5}{The architecture in which both instructions and data are held in the same unified memory space is named after:}{Blaise Pascal}{Charles Babbage}{John von Neumann}{Linus Torvalds}

\mcqitem{6}{Which historical mechanical calculating device was invented by Blaise Pascal in 1642 to perform addition and subtraction?}{Analytical Engine}{Pascaline}{ENIAC}{Abacus}

\mcqitem{7}{Who designed the mechanical programmable "Analytical Engine" in the 1830s containing concepts of memory and processor?}{Herman Hollerith}{Charles Babbage}{John von Neumann}{Alan Turing}

\mcqitem{8}{Herman Hollerith successfully used which technology to automate data processing for the 1890 US Census?}{Vacuum tubes}{Punched cards}{Transistors}{Integrated circuits}

\mcqitem{9}{What was the primary electronic switching component used in first-generation computers (1940--1956)?}{Transistors}{Vacuum tubes}{Integrated circuits}{Microprocessors}

\mcqitem{10}{Which famous early computer is a classic example of a first-generation electronic computer?}{Apple II}{ENIAC}{IBM PC}{Cray-1}

\mcqitem{11}{Second-generation computers (1956--1963) replaced vacuum tubes with:}{Transistors}{Microprocessors}{Silicon wafers}{Mechanical gears}

\mcqitem{12}{Which programming languages became prominent during the second generation of computers?}{Java and Python}{FORTRAN and COBOL}{HTML and CSS}{C\# and Swift}

\mcqitem{13}{Third-generation computers (1964--1971) were characterized by the widespread adoption of:}{Punched cards}{Integrated Circuits (ICs)}{Vacuum tubes}{Artificial intelligence}

\mcqitem{14}{Which generation of computers introduced keyboards, monitors, and modern operating systems to replace batch punch cards?}{First generation}{Second generation}{Third generation}{Fifth generation}

\mcqitem{15}{The fourth generation of computers (from 1971 onward) is fundamentally defined by the development of:}{Vacuum tubes}{Discrete relays}{Microprocessors}{Optical gears}

\mcqitem{16}{Which company introduced the world's first commercial single-chip microprocessor (the 4004) in 1971?}{IBM}{Intel}{Apple}{Microsoft}

\mcqitem{17}{Fifth-generation computing focuses primarily on which technological direction?}{Vacuum tubes}{Artificial Intelligence and parallel processing}{Pure mechanical gears}{Magnetic drums}

\mcqitem{18}{The computer characteristic stating that "incorrect input data inevitably produces incorrect output information" is known as:}{FIFO}{GIGO (Garbage In, Garbage Out)}{LIFO}{DMA}

\mcqitem{19}{A computer's ability to perform repetitive operations continuously for hours without fatigue or loss of concentration is called:}{Versatility}{Speed}{Diligence}{Automation}

\mcqitem{20}{Which characteristic refers to a computer's capability to perform widely diverse tasks ranging from music playback to complex calculations?}{Accuracy}{Versatility}{Diligence}{Volatility}

\mcqitem{21}{Which CPU component is responsible for performing additions, subtractions, and logical comparisons?}{Control Unit}{Arithmetic Logic Unit (ALU)}{Instruction Register}{Address Bus}

\mcqitem{22}{Which component acts as the supervisor of the CPU, generating timing and control signals to direct operations?}{Control Unit (CU)}{ALU}{Cache}{RAM}

\mcqitem{23}{The high-speed system bus that carries the actual data and instructions between CPU, memory, and devices is the:}{Address Bus}{Control Bus}{Data Bus}{Power Bus}

\mcqitem{24}{Which system bus is unidirectional and transmits memory locations from the CPU to primary memory?}{Data Bus}{Address Bus}{Control Bus}{Expansion Bus}

\mcqitem{25}{Which system bus carries command signals such as Memory Read, Memory Write, and Interrupt acknowledgments?}{Control Bus}{Data Bus}{Address Bus}{Local Bus}

\mcqitem{26}{Primary memory directly accessed by the CPU includes:}{HDD and SSD}{RAM and Cache}{Magnetic tape and Optical disks}{USB drives}

\mcqitem{27}{Which memory loses all its contents immediately when electrical power is switched off?}{Read-Only Memory (ROM)}{Random-Access Memory (RAM)}{Solid-State Drive (SSD)}{Hard Disk Drive (HDD)}

\mcqitem{28}{What is the primary function of Read-Only Memory (ROM)?}{To store temporary variables during gaming}{To hold non-volatile startup firmware (e.g., BIOS)}{To speed up web browsing}{To back up user photos}

\mcqitem{29}{The physical, tangible parts of a computer system that you can touch are collectively known as:}{Software}{Hardware}{Firmware}{Freeware}

\mcqitem{30}{Programs and sets of instructions that direct physical computer hardware to perform tasks are called:}{Hardware}{Software}{Peripherals}{Chipsets}

\mcqitem{31}{An operating system such as Windows, Linux, or macOS is an example of:}{Application software}{System software}{Utility software}{Programming firmware}

\mcqitem{32}{Software designed to perform specific user-oriented tasks, such as word processing or web browsing, is called:}{System software}{Application software}{Device driver}{Operating system}

\mcqitem{33}{Small software programs that allow an operating system to communicate with a specific hardware peripheral (like a printer) are:}{Device drivers}{Spreadsheets}{Compilers}{Firmware updates}

\mcqitem{34}{Low-level software permanently written into non-volatile ROM chips to initialize hardware during startup is known as:}{Application software}{Firmware}{Freeware}{Shareware}

\mcqitem{35}{Which type of computer is designed for high-performance, complex scientific modeling, weather forecasting, and molecular research?}{Microcontroller}{Supercomputer}{Personal computer}{Embedded computer}

\mcqitem{36}{Very large, powerful, and reliable computers used by banks and airlines for massive multi-user transaction processing are:}{Workstations}{Mainframe computers}{Laptops}{Analog computers}

\mcqitem{37}{A computer dedicated to providing shared resources, files, databases, or web pages to client computers across a network is a:}{Client}{Server}{Transistor}{Peripheral}

\mcqitem{38}{A single-user, high-performance desktop computer equipped with powerful graphics and CPU for engineering and CAD design is a:}{Mainframe}{Workstation}{Embedded system}{Supercomputer}

\mcqitem{39}{A small computer built directly into an everyday appliance (such as an automobile ABS or washing machine) to perform a dedicated task is a/an:}{Supercomputer}{Embedded computer}{Mainframe}{Server}

\mcqitem{40}{An integrated circuit that combines a processor core, memory, and input/output interfaces onto a single silicon chip is a:}{Microcontroller}{Supercomputer cluster}{Mainframe rack}{Graphics card}

\mcqitem{41}{Computers that process information represented by continuously variable physical quantities, such as voltage or pressure, are:}{Digital computers}{Analog computers}{Hybrid computers}{Mainframe computers}

\mcqitem{42}{Modern personal computers and laptops that process discrete binary digits (0s and 1s) are classified as:}{Analog computers}{Digital computers}{Continuous computers}{Mechanical engines}

\mcqitem{43}{A computer that combines the measurement capabilities of analog systems with the logical precision of digital systems is called a:}{Supercomputer}{Hybrid computer}{Microcontroller}{Mainframe}

\mcqitem{44}{A computer system designed to run a wide variety of software applications (e.g., word processor, compiler, games) is called:}{Special-purpose computer}{General-purpose computer}{Dedicated controller}{Analog meter}

\mcqitem{45}{An embedded controller designed solely to operate the braking system in an automobile is an example of a:}{General-purpose computer}{Special-purpose computer}{Mainframe computer}{Supercomputer}

\mcqitem{46}{How many bits are in one standard Byte?}{2}{4}{8}{16}

\mcqitem{47}{Which unit represents approximately one trillion bytes in decimal metric measurement ($10^{12}$ bytes)?}{1 Gigabyte}{1 Terabyte}{1 Megabyte}{1 Kilobyte}

\mcqitem{48}{Processor clock speed is commonly measured in:}{Bytes per second (Bps)}{Gigahertz (GHz)}{Dots per inch (DPI)}{Frames per second (FPS)}

\mcqitem{49}{What is a fundamental limitation of modern digital computers?}{They lack high processing speed}{They cannot store large quantities of data}{They lack human judgment, common sense, and independent moral reasoning}{They cannot execute repetitive tasks}

\mcqitem{50}{Which component converts alternating current (AC) electricity from the wall socket into regulated low-voltage direct current (DC) for computer parts?}{Motherboard}{Power Supply Unit (PSU)}{Heat sink}{Central Processing Unit (CPU)}

\newpage
\section{Answer Key \& Explanations --- Unit 2}

\begin{answerkeytable}{Unit 2: Computer Fundamentals --- Answer Key}
\begin{tabular}{*{10}{c}}
\toprule
\textbf{Q} & \textbf{Ans} & \textbf{Q} & \textbf{Ans} & \textbf{Q} & \textbf{Ans} & \textbf{Q} & \textbf{Ans} & \textbf{Q} & \textbf{Ans} \\
\midrule
1 & \textbf{B} & 11 & \textbf{A} & 21 & \textbf{B} & 31 & \textbf{B} & 41 & \textbf{B} \\
2 & \textbf{B} & 12 & \textbf{B} & 22 & \textbf{A} & 32 & \textbf{B} & 42 & \textbf{B} \\
3 & \textbf{B} & 13 & \textbf{B} & 23 & \textbf{C} & 33 & \textbf{A} & 43 & \textbf{B} \\
4 & \textbf{B} & 14 & \textbf{C} & 24 & \textbf{B} & 34 & \textbf{B} & 44 & \textbf{B} \\
5 & \textbf{C} & 15 & \textbf{C} & 25 & \textbf{A} & 35 & \textbf{B} & 45 & \textbf{B} \\
6 & \textbf{B} & 16 & \textbf{B} & 26 & \textbf{B} & 36 & \textbf{B} & 46 & \textbf{C} \\
7 & \textbf{B} & 17 & \textbf{B} & 27 & \textbf{B} & 37 & \textbf{B} & 47 & \textbf{B} \\
8 & \textbf{B} & 18 & \textbf{B} & 28 & \textbf{B} & 38 & \textbf{B} & 48 & \textbf{B} \\
9 & \textbf{B} & 19 & \textbf{C} & 29 & \textbf{B} & 39 & \textbf{B} & 49 & \textbf{C} \\
10 & \textbf{B} & 20 & \textbf{B} & 30 & \textbf{B} & 40 & \textbf{A} & 50 & \textbf{B} \\
\bottomrule
\end{tabular}
\end{answerkeytable}

\begin{expbox}[Selected Question Explanations]
\begin{itemize}[leftmargin=5mm]
    \item \textbf{Q.5 (C):} The von Neumann architecture introduced the stored-program concept in 1945 where code and data reside in common memory.
    \item \textbf{Q.9 (B):} First-generation computers (e.g., ENIAC) used fragile, heat-producing vacuum tubes as electronic switches.
    \item \textbf{Q.11 (A):} Transistors made second-generation computers much smaller, faster, and more reliable than vacuum tubes.
    \item \textbf{Q.16 (B):} The Intel 4004 (1971) was the first commercial single-chip microprocessor, kicking off the 4th generation.
    \item \textbf{Q.18 (B):} GIGO (Garbage In, Garbage Out) means that the accuracy of a computer's output depends strictly on the correctness of input data and software.
    \item \textbf{Q.24 (B):} The Address Bus is unidirectional; it carries memory addresses from the CPU to memory devices.
    \item \textbf{Q.39 (B):} Embedded computers are dedicated microcontrollers placed inside washing machines, cars, or medical devices to carry out specific control routines.
    \item \textbf{Q.41 (B):} Analog computers represent data as continuous physical variables (voltage, temperature), whereas digital computers use discrete bits ($0, 1$).
\end{itemize}
\end{expbox}

\newpage
